\chapter{DNA storage and error-correcting design}
A tube does not contain a file in byte order. It contains physical copies of strands,
some missing, some changed, and some sampled many more times than others. The reader
must infer which fragments belong together, recover absent fragments when the code
permits it, and decide whether the resulting bytes are trustworthy. Storage is thus
an end-to-end communication problem, not merely a translation from four letters
to two bits.

Our worked target is explicit: recover the 28-byte message
\texttt{DNA stores bytes, not proof.} from an unordered digital strand pool after
one located erasure in each parity group. We then derive a separate constrained
sequence code. The separation is intentional: the first reference proves
framing and recovery; the second proves membership and reversible ranking under
two sequence constraints. Neither is a synthesis-ready oligonucleotide design.

\section{Define the channel before choosing the code}
\index{DNA storage!channel}
Let a source file be $u\in\{0,1\}^B$. An encoder produces intended
sequences $S=\{s_1,\ldots,s_M\}$. Writing creates a molecular population with copy
counts $n_i$ and potentially changed sequences. Storage and retrieval alter that
population. Sequencing samples reads $R$, and a decoder maps them to a file
estimate $\widehat u$ or declared failure:
\begin{equation}
u \xrightarrow{\mathrm{encode}} S
\xrightarrow{\mathrm{write/store/access/read}} R
\xrightarrow{\mathrm{decode}} \widehat u\ \text{or failure}.
\label{dnad25:channel}
\end{equation}
A decoder that always returns plausible text can hide uncertainty. A decoder
that fails when evidence is insufficient makes the claim testable.

\Plate{01-channel}{A digital codeword collection becomes a population of molecules,
then a biased sample of reads. The numbered boundaries distinguish addressing,
physical copy counts and reconstruction. The reference executes the digital
boundaries, not the biochemical ones.}{dnad25:channel-plate}

A substitution changes a base while retaining position; an insertion or deletion
shifts alignment. A strand erasure removes an entire intended codeword from usable
evidence. A file error means the reconstructed source differs anywhere. A small
per-base error rate does not guarantee file recovery if an indispensable strand
is missing. Conversely, noisy reads may reconstruct a file exactly when their
errors and coverage fit the decoder's assumptions.

Copy-number skew is not simply another spelling of base error. Two strands may both
be error-free while one appears in a thousand reads and the other in none. Research
on molecular bias distinguishes synthesis and amplification effects;
its reported models are context, not calibration for our digital channel
\cite{bias}. Chapter 24's occupancy and capture calculations explain why nominal
inventory and observed coverage cannot be interchanged.

\section{The two-bit ceiling and the rate actually delivered}
\index{coding rate}
Four equiprobable unconstrained symbols can carry $\log_2 4=2$ bits per nucleotide.
This is an alphabet ceiling, not an achieved archive density. An intended oligo of
length $L$ may devote bases to physical access, a digital index and redundancy.
Sequence restrictions further reduce the usable words. Distinguish
\begin{equation}
R_{\rm logical}=\frac{B}{\sum_{i=1}^{M}|s_i|},
\qquad
R_{\rm physical}=\frac{B}{\sum_{i=1}^{M} n_i|s_i|}.
\label{dnad25:rate}
\end{equation}
Both denominators count nucleotides, but the second counts all molecule copies.
It must use a declared inventory, not a sequencing read count. A mass-density
comparison additionally needs a molecular-mass and double/single-strand
convention, container/encoding overhead and measured material quantities.

Our byte map reads the most significant bit pairs first:
$00\mapsto A$, $01\mapsto C$, $10\mapsto G$, $11\mapsto T$.
Byte 173 is binary 10101101, so its word is \texttt{GGTC}.
The tests independently recover the base-four digits of every possible byte.
A byte occupies four nucleotides; an odd-length arbitrary bit string would need
a separate length/padding contract. This reference accepts bytes, not a silently
padded bit stream.

\section{An outer erasure code derived from XOR}
\index{erasure code}
Split a source segment into equal-width byte vectors $A,B\in\{0,1\}^{8w}$ and
store three shards:
\begin{equation}
(A,B,P),\qquad P=A\oplus B.
\label{dnad25:parity}
\end{equation}
If $A$ is missing, $P\oplus B=A$; if $B$ is missing, $P\oplus A=B$.
If only parity is missing, $A,B$ already hold the data. Any two known,
correct, distinct-address shards reconstruct the segment. With two erasures,
one known vector cannot determine two independent unknown vectors.

This is a $(3,2)$ single-parity code over bytewise XOR. Its parity-check relation
is $A\oplus B\oplus P=0$. In each bit coordinate the minimum Hamming distance
is two: changing one data bit also changes its parity bit. It locates no arbitrary
error by itself. If all three values arrive but violate the relation, the reference
rejects them instead of choosing one to trust.

\Plate{04-redundancy}{One parity equation reconstructs one located erasure.
All three present-but-inconsistent shards trigger failure; two missing shards
are underdetermined. Addresses and validity checks precede this algebra.}
{dnad25:redundancy}
\Listing{recover_group}

A common implementation defect is \texttt{zip(a,b)} on unequal byte strings:
Python stops at the shorter input. A parity calculation would silently
discard a suffix. The reference validates equal widths before XOR, and tests
this failure explicitly. Mathematical notation assumes compatible vector
spaces; software must enforce that assumption.

\section{Frame the payload, then decode an unordered pool}
\index{DNA storage!framing}
Each shard carries a 12-byte header: version (one byte), group index (two),
slot (one), chunk width (two), original file length (four) and group count (two).
The payload contributes $w$ bytes; CRC-32 contributes four. Integers use network
big-endian order. Version 1, $0\leq\mathrm{slot}\leq2$, and the manifest's
layout fields must agree. The reference limits $1\leq w\leq4096$ and at most
65,535 groups to make its finite implementation contract explicit.

The last source pair is zero-padded, but original byte length is retained.
Without it, a real trailing zero byte is indistinguishable from padding.
The archive also has an out-of-band manifest containing a SHA-256 digest of the
source. Header redundancy makes reads interpretable; the manifest
fixes which layout and final source are expected.

\Listing{encode_archive}

Decoding validates the manifest, converts canonical DNA to bytes, checks
frame length and CRC, and verifies address/layout. Invalid reads become
unusable evidence. Exact duplicate valid reads are counted once at their address;
different CRC-valid payloads at one address trigger failure. Groups are reconstructed
in index order and the concatenation is trimmed to original length. Finally the
file digest must agree.

\Listing{decode_archive}

The worked archive has four groups and 12 unique strands, each 80 nt long:
$(12+4+4)$ bytes times four nucleotides per byte. Its net logical rate is
$224/960=0.233333$ bits per unique nucleotide. This low rate reflects deliberately
large teaching headers and parity, not a limit on production storage systems.
Delete slot $A$ in every group, reverse the remaining pool order, and repeat one
read. The decoder recovers the exact 28 bytes, reports four reconstructed groups,
eight unique addresses and one duplicate. This is a digital intervention,
not a dropout measurement.

\section{Checksums detect errors; they do not authenticate a story}
\index{checksum}\index{authentication}
CRC-32 is an accidental-error check, not an adversarial integrity guarantee.
A changed payload with a recomputed CRC is still a valid frame. Parity also allows
coherent alternative pairs: replacing $A,B$ and recomputing $P$ preserves
the relation. The final digest catches the deliberately corrupted file
when compared with the original trusted manifest.

If an attacker can replace both data and manifest, the digest is replaced too.
Authenticating the expected manifest requires an external trust mechanism.
We make no universal $2^{-32}$ error-probability assertion for an arbitrary
correlated DNA channel: that probability needs a distribution over corruptions
and checksums. Our exhaustive single-base tests establish detection on their
finite cases, not a theorem covering every possible channel.

The tests replace each position in a worked strand by each other canonical
base. Every changed frame is rejected and the remaining two shards recover
the data. They also exercise length-changing reads, all three erasure locations,
all small byte pairs, conflicting valid duplicates and coherent parity-valid
corruption detected by the file digest. These checks address different mechanisms.

\section{Consensus requires a declared alignment model}
\index{consensus!alignment}
For an already identified, equal-length, aligned cluster, count A/C/G/T at each
column. Return a base only if its count exceeds half the cluster size.
If any column lacks a strict majority, return an erasure. A lexical tie-break
turns ignorance into a systematic A preference; it is not scientific evidence.

\Plate{02-errors}{A substitution preserves a coordinate; an insertion shifts it.
Column consensus applies only after identity, orientation and alignment are
established. A strict-majority tie is routed to erasure, not invented sequence.}
{dnad25:errors}
\Listing{strict_consensus}

The finite example generates 11 reads of one strand with an assigned independent
substitution probability of 0.02 and seeds 0 through 10. Their consensus matches
the intended strand. This is one simulated realization, not an
estimated channel success rate. Consensus cannot repair a missing entire cluster,
and repeated systematic errors can survive a majority vote.

An indel-aware decoder must infer alignment, distinguish strand identities,
use quality information where calibrated, and propagate ambiguous reconstruction.
Our equal-length routine rejects length differences. The framing
routine treats unusable indel reads as erasures, rather than pretending it has
performed alignment. Edit-error correction is outside this implementation.

\section{Count constrained words instead of hoping a screen succeeds}
\index{constrained coding}
The raw archive proves recovery but its first strand has GC fraction 0.225
and a homopolymer run of 18. It fails the constraints we derive next.
Discarding that strand without retaining an invertible mapping loses source
information. A constrained encoder maps messages into its legal language
by construction.

Let $\mathcal C(n,g_{\min},g_{\max},r)$ contain words of length $n$,
GC count in $[g_{\min},g_{\max}]$, and no repeated-base run longer than $r$.
Define $C(\ell,a,j,g)$ as the number of legal completions with $\ell$ positions
remaining, last base $a$, current run $j$ and GC count $g$. At the terminal state,
\[
C(0,a,j,g)=\mathbf 1\{g_{\min}\leq g\leq g_{\max}\}.
\]
For $\ell>0$,
\begin{equation}
C(\ell,a,j,g)=
\sum_{\substack{b\in\{A,C,G,T\}\\j'(a,j,b)\leq r}}
C(\ell-1,b,j'(a,j,b),g+\mathbf 1\{b\in\{G,C\}\}),
\label{dnad25:count}
\end{equation}
where $j'=j+1$ if $a=b$ and $j'=1$ otherwise.
States with $g>g_{\max}$ or $g+\ell<g_{\min}$ contribute zero.
The suffix count depends on a small sufficient state, not every spelling
of the prefix. Memoization shares subproblems; naive enumeration inspects
$4^n$ candidates.

\Plate{03-constraints}{Legal completion counts partition the lexical codebook
into intervals. Ranking and unranking use those interval sizes; a rejected
homopolymer branch has zero completions. Rank 173 maps to AACGCTGT.}
{dnad25:constraints}
\Listing{ConstraintCodec.__init__}

The finite teaching domain is $n\leq64$. Counting uses exact integers,
but run bounds and GC windows still determine cache memory and work.
This is not a claim that arbitrary biochemical constraints remain tractable
with the same sufficient state.

\section{Rank and unrank: a reversible encoder}
Order words lexically with $A<C<G<T$. For rank $q$, inspect next bases in
that order. If a branch has $c$ completions and $q\geq c$,
skip its interval and subtract $c$. Otherwise choose that base and continue.
Each choice preserves $0\leq q<$ the selected suffix count.
This invariant proves that the terminal word has the requested rank.

The inverse adds completion counts for all smaller branches at each position.
Thus rank and unrank use the same partition in opposite directions.
Tests also compare both with an independently enumerated, sorted legal
codebook, avoiding a shared-bug roundtrip alone.

\Listing{ConstraintCodec.unrank}
\Listing{ConstraintCodec.rank}

For $n=8$, exactly four GC bases and $r=2$, there are 14,696 words.
A fixed-length binary block uses
$k=\lfloor\log_2|\mathcal C|\rfloor=13$ bits, giving $1.625$ bits/nt.
Ranks 0 through 8,191 leave 6,504 legal words unused.
This rate gap pays for constraints and a power-of-two fixed block;
it is not automatically the channel's Shannon capacity.
The reference supports legal ranks beyond the binary subset.

\par\medskip\noindent\input{\Generated/code_sizes.tex}\par\medskip
Concatenation requires its own boundary contract: two run-limited words
can form a long homopolymer at their join. GC bounds on payload words do not
constrain appended headers or primers. The constrained code is a separate
mechanism, not silently grafted onto the raw archive.

\section{Random access is a molecular selection problem}
\index{DNA storage!random access}
A digital index tells a decoder where a read belongs after observation.
A physical address selects which molecules will be observed.
Primer-mediated amplification is one strategy; specificity depends
on target and off-target interactions, not the presence of a unique
integer in a header.

\Plate{05-access}{A selected primer pair enriches a target-address family.
Off-target interactions and unequal amplification change the read population.
This is a conceptual mechanism, not a sequence-validated primer design.}
{dnad25:access}

Organick et al. report random access to individual files in a large
encoded collection \cite{organick}. This is a demonstration under
their system, not proof that our addresses are biochemically orthogonal.
The Goldman et al. abstract reports exact reconstruction of
739 kilobytes \cite{goldman}; it cannot supply our missing synthesis,
sequencing and independent experimental validation.

If a target strand's usable reads follow Poisson mean $\lambda_i$,
its erasure probability is $e^{-\lambda_i}$.
Uniform means give Chapter 24's occupancy model; heterogeneous means
show why average coverage can hide dropout. State independence and
Poisson assumptions before multiplying strand success probabilities.

\section{A rate--recovery frontier, not a magic storage claim}
\Plate{06-boundary}{Storage preserves an intended message across a noisy channel;
molecular computation transforms a state under a mechanism.
They share coding, assembly and readout but have different success criteria
and resource denominators.}{dnad25:boundary}

Redundancy can improve recovery for an assumed channel but requires more unique
sequence and potentially more material. Physical copies can reduce sampling
erasure while enlarging the inventory denominator. Sequence restrictions can
reduce particular unwanted interactions while lowering code rate. These
are coupled engineering choices, not independent magic improvements.

A serious evaluation reports source size, unique strand count/length,
logical rate, physical copy distribution, write/storage/read conditions,
error model, decoder version, access method, failed attempts and file verification.
The reference supplies inspectable submechanisms for reasoning about that
system, not a substitute for physical characterization.

\section{Connect the finite mechanisms to current research}
The purpose of a small reference is to make a larger system intelligible,
not to equate its performance with that system. Gimpel et al.'s 2023 digital
twin explicitly treats workflow-dependent error and coverage bias
\cite{twin}. Its scope reminds us that a single substitution probability
cannot characterize every synthesis, amplification and aging condition.
Transfer to a different process requires new evidence.

Bar-Lev et al.'s 2025 study combines neural reconstruction with error-correcting
codes and a safety-margin mechanism \cite{deep-storage}.
Our equal-column vote implements none of that pipeline; learned reconstruction
also needs a declared training-channel distribution and out-of-distribution
failure tests. The paper's later correction adds missing prior work on
simulated-data generation \cite{deep-correction}. A professional literature
ledger tracks such corrections, not only attractive performance summaries.

Li et al.'s 2026 Siyuan Code study joins a forbidden-pattern automaton,
dynamic-programming counts and bijective coding \cite{bijective}.
Our two-constraint suffix counter teaches the general counting/ranking idea;
it does not reproduce their automaton, full-strand codec, experimental density
or production timing. Their distinction between coding potential, net rate
and physical density reinforces Eq.~\ref{dnad25:rate}.
General enumerative coding is established theory, not a novelty claim here.

These approaches can address different bottlenecks: a channel model guides
redundancy, a reconstruction method interprets noisy observations, and a
constrained codec selects usable words. Combining them requires matched
sequence length, physical copy count, channel conditions, source workload
and file-level failure criteria. A reported speed ratio under one workload
does not imply that the complete write--store--read system is faster by
that factor.

\section{Twenty worked exercises}
\begin{enumerate}
\item Encode byte 173. \textbf{Answer:} 10101101 gives GGTC.
\item Why retain source length? \textbf{Answer:} padding and genuine trailing
zero bytes otherwise cannot be distinguished.
\item Derive missing $A$. \textbf{Answer:} $(A\oplus B)\oplus B=A$.
\item Can one shard recover two unknown shards? \textbf{Answer:} no;
the parity relation admits multiple source pairs.
\item Why validate XOR widths? \textbf{Answer:} a truncated zip omits coordinates.
\item For $w=4$, find strand length. \textbf{Answer:} $4(12+4+4)=80$ nt.
\item Find net rate for 28 bytes in 12 such strands. \textbf{Answer:}
$8(28)/(12\cdot80)=0.233333$ bits/nt.
\item Double all physical copies. \textbf{Answer:} logical rate stays fixed;
physical inventory rate halves.
\item What does a duplicate valid read add? \textbf{Answer:} read-count evidence,
not another independent parity equation at that address.
\item Conflicting CRC-valid duplicates? \textbf{Answer:} fail, not last-write-wins.
\item Does CRC authenticate? \textbf{Answer:} no; an adversary can recompute it.
\item Why can coherent corruption pass parity? \textbf{Answer:}
$P'=A'\oplus B'$ satisfies the relation for changed data.
\item Why a trusted manifest? \textbf{Answer:} replacing file and expected
digest defeats comparison against that altered digest.
\item A versus C in two reads? \textbf{Answer:} strict consensus returns erasure.
\item Can column voting repair insertion automatically? \textbf{Answer:} no;
the coordinate/alignment model changed.
\item Derive GC pruning. \textbf{Answer:} $g>g_{\max}$ violates the upper bound;
$g+\ell<g_{\min}$ cannot reach the lower bound.
\item Unranking invariant? \textbf{Answer:} residual rank remains in the chosen
branch's completion interval.
\item Capacity for 14,696 words? \textbf{Answer:} 13 bits since
$8192\leq14696<16384$.
\item Why is concatenation not run-safe? \textbf{Answer:} runs join across a boundary.
\item Design an independent validation ladder. \textbf{Rubric:}
all-byte oracle; exhaustive short codebooks; rank trace; all single erasures;
malformed/conflicting reads; alignment failures; coherent corruption;
unordered pools and padding; declared physical channel; write/read
experiments with independent preparations and trusted source verification.
\end{enumerate}

\section{Bridge: from decoder correctness to experiment evidence}
Chapter 26 turns these boundaries into a laboratory provenance contract.
A decoder test proves a statement about a digital function. Demonstrating
that a molecular pool meets its assumptions requires separate material,
preparation, measurement and analysis records. Keeping these identities
separate makes an elegant code an auditable experimental claim.
